\chapter{Speed and memory}
\label{chap:performance}

I treated speed and memory as constraints from the start, but never at the expense of correctness:
each optimisation came after the result it served was right, and I checked that it left the output
unchanged (byte-identical, or with every difference explained). Today the whole stack is segmented in
about \pbValueSeconds\,s with a peak of \pbValuePeakGB\,GB of memory.

\paragraph{Test machine} Timings depend on the hardware, so they are only meaningful next to it. Every
timing of this report (here and in Table~\ref{tab:history}) was measured on the same machine, recorded
automatically when the runs were timed:
\begin{center}
\small
\begin{tabular}{@{}ll@{}}
\toprule
Processor        & \pbValueCpu{} (\pbValueCores{} logical cores) \\
Memory           & \pbValueMemoryGB\,GB \\
System           & \pbValueSystem \\
Power profile    & \pbValuePowerProfile{} (every version timed in one session) \\
Python           & \pbValuePython{} (report tool) \\
pyBones threads  & \pbValueWorkers{} (\code{pipeline.workers}) \\
\bottomrule
\end{tabular}
\end{center}
The times are wall-clock times of the whole command, including the start-up of each version's
environment; the memory is the peak resident size of the process tree. On another machine the absolute
values will differ; the comparison between versions holds, since they all ran on this one.

\paragraph{What makes it fast}
\begin{itemize}
  \item \textbf{Vectorisation}: no Python loop ever runs over pixels or labels. Component statistics are
        one \code{bincount}, morphological reconstruction one labelling, the Otsu search one array
        expression over all threshold pairs, the eigenvalues a closed formula over millions of matrices.
  \item \textbf{OpenCV} for every filter: Gaussian blurs, box filters (integral images), grey-level
        dilations and erosions, connected components, separable Hessian filters.
  \item \textbf{Threads} for the per-slice work (about twice as fast with four threads).
  \item \textbf{Work only where it matters}: the Hessian eigenvalues are computed only on the voxels that
        can become ambiguous, a few percent of each slice instead of a quarter of it.
\end{itemize}

\paragraph{What keeps memory low}
\begin{itemize}
  \item \textbf{Streaming}: slices go through bounded windows; volumes are written slice by slice and
        read back memory-mapped. A full 32-bit label volume (1.4\,GB) or six float Hessian volumes
        (8\,GB) are never built.
  \item \textbf{Compact buffers}: the slices waiting for the 3D burst box are bit-packed (eight times
        smaller), the burst density lives on a 12\um{} grid (sixteen times smaller), the tubularity is
        stored on one byte.
  \item \textbf{Chunks}: eigenvalues, mesh smoothing, normals and mesh writing process bounded
        batches.
\end{itemize}

\paragraph{Meshes} Extracting and smoothing the surface of the vessels takes
\pbValueMeshVesselSeconds\,s and that of the bone \pbValueMeshBoneSeconds\,s on the same machine, most
of it in the Taubin smoothing; writing the file adds a few seconds. A mesh job reports
four weighted stages (block occupancy, surface, smoothing, file), so its progress bar moves with the
work, and it can be cancelled between them.

\paragraph{Progress} Every long operation reports weighted stages, and its progress bar and estimate
of the time left are only right if each weight matches the stage's share of the run. Weights set once
drift as the code gets faster in one place and not another, so the report tool measures them on the
sample every time it runs: a stopwatch listens to each stage's progress reports, and a stage lasts
from the previous stage's last report to its own last one, so the work before its first report counts
too. Table~\ref{tab:stage-weights} compares each weight with its measured share; the tool warns when
one is more than \pbValueStageWeightTolerance{} off, and the weights of this version were set from these
measurements. Some stages are one long step that cannot report its own progress (the skeleton, the
surface of a mesh): there, the page keeps the time ticking, and the server sends every new stage at
once, so that a long silent stage is never shown under the name of the previous one.

\begin{table}[ht]
  \centering
  \small
  \begin{tabular}{llrr}
\toprule
Operation & Stage & Weight & Measured share \\
\midrule
segmentation & thresholds & 0.020 & 0.022 \\
 & classify & 0.650 & 0.631 \\
 & analyse & 0.110 & 0.114 \\
 & resolve & 0.170 & 0.179 \\
 & clean & 0.050 & 0.055 \\
mesh & occupancy & 0.300 & 0.296 \\
 & surface & 0.090 & 0.094 \\
 & smoothing & 0.530 & 0.532 \\
 & writing & 0.080 & 0.077 \\
quantification & cavities & 0.350 & 0.345 \\
 & skeleton & 0.580 & 0.583 \\
 & calibre & 0.050 & 0.055 \\
 & network & 0.020 & 0.017 \\
\bottomrule
\end{tabular}

  \caption{Progress stages of the three long operations on the sample: configured weight (normalised
    over the stages that run) and measured share of the run's time.}
  \label{tab:stage-weights}
\end{table}

\paragraph{On the web page} A page that stays open for hours must not grow: an image cache that never
empties is enough to crash a browser. So the page holds every resource once:
\begin{itemize}
  \item \textbf{Slices are never kept.} Each one is fetched with the browser's cache disabled, shown
        through a temporary \code{blob:} address that is released as soon as the image is decoded, and
        the server forbids caching it as well. Scrolling through a stack for hours costs one image.
  \item \textbf{One slice request at a time.} Dragging the slider never piles renders up on the
        server: a position asked for during a render waits, and only the latest one is fetched.
  \item \textbf{Meshes live on the graphics card.} The CPU copy that three.js keeps after uploading a
        mesh is released, which halves the browser's share; a segmentation left behind frees its
        meshes. A mesh may be kept by the browser between visits, but it is checked each time, so a
        rebuilt mesh is never hidden by an old copy.
  \item \textbf{Following a job holds no thread.} Job events reach the browser through the server's
        event loop, so any number of open pages cannot starve the slice and list requests.
\end{itemize}

My last optimisation pass brought the segmentation from 36\,s and 3.3\,GB to the figures above, and the
smoothing of the bone mesh (about 7 million triangles) from 1.9\,GB to 1.2\,GB with an identical result.
Version 0.1 was faster and lighter still, but it computed much less: the 3D refinement and the
tubularity are what the time and memory now buy.
